\part{Reinforcement learning and intelligent scheduling}

\begin{warnbox}{The RL result in one sentence}
Learned priority policies (CEM and PPO) were trained, integrated into the live loop, and evaluated against the
rule-based demand-aware policy. \textbf{On the network they give no significant gain}: in emulation, PPO vs proposed
$\Delta U=\NPpoDelta$ ($p=\NPpoP$). In the calibrated simulator, reward per step is \NSeqRuleReward{} (rule) vs
\NSeqCemReward{} (CEM) vs \NSeqPpoReward{} (PPO). That gain is statistically significant on paired test seeds but
tiny (+0.5\,\% and +0.3\,\%), and comes with slightly more SLA violations per episode
(\NSeqRuleViol{} vs \NSeqCemViol{} vs \NSeqPpoViol). Every live decision is logged as an RL transition, and an online-CEM
variant learns from verified live transitions, but it shows no measurable trend over 10 runs. Present RL as an
evaluated alternative with a negative live result, not as a contribution.
\end{warnbox}

\section{Does the RL component solve a real problem?}

\begin{qa}{541}{What specific decision is the RL scheduler meant to learn?}
Per-context (CEM) or per-state (PPO) multipliers $m_s\in[0.1,10]$ on the four slices' priority weights. The allocator
then computes rates. RL never sets rates directly.
\end{qa}

\begin{qa}{542}{Why is reinforcement learning appropriate for this decision?}
The allocator maximises weighted max-min fairness, not the evaluation objective $U$, and allocations have delayed
effects (backlog, TCP adaptation, outages). RL can in principle optimise $U$ end to end over those dynamics.
Whether it \emph{helps} is an empirical question, and our answer is no.
\end{qa}

\begin{qa}{543}{Why not a deterministic optimiser or a well-designed heuristic?}
That is the baseline we use, and it wins (ties). For a single link with known demands, a convex program for $U$ would
be the principled alternative to RL. Its gain would likely be equally small, because HTB borrowing redistributes
unused capacity in the data plane anyway.
\end{qa}

\begin{qa}{544}{What does the RL agent observe?}
PPO (21 dimensions): context one-hot (8), demand/$C$ (4), normalised backlog (4), current allocation/$C$ (4), and an
on-backup flag. Live, the same vector is built from the controller state and HTB backlog (\code{observation_from_live}).
CEM observes only the context.
\end{qa}

\begin{qa}{545}{What actions can it take?}
A continuous $a\in[-1,1]^4$, mapped to multipliers $10^{a_s}$.
\end{qa}

\begin{qa}{546}{What is the reward function?}
Per step: $r=100\sum_s \tilde w_s\,\text{sat}_s-10\cdot\#\{s:\text{delay}_s>\text{SLA}_s\}-0.5\cdot\overline{\text{loss}\%}$,
where $\tilde w$ are the \emph{context's} normalised weights (not the agent-scaled ones, which would let it inflate its
own reward), and sat $=\min(1,\text{served}/\text{demand})$.
\end{qa}

\begin{qa}{547}{How does the reward balance SLA compliance, throughput, fairness and reconfiguration cost?}
Throughput through satisfaction; SLA through the violation penalty; loss through a penalty.
\honest{There is no fairness term and no reconfiguration cost. Fairness is evaluated (Jain) but not rewarded.}
\end{qa}

\begin{qa}{548}{How do you prevent the agent maximising one metric while harming critical services?}
With a safety layer outside the policy:
\begin{itemize}
\item multipliers are clipped;
\item floors are allocated before the policy's shares;
\item healing constraints apply after;
\item reward weights are the context's, not the agent's.
\end{itemize}
The policy can reshuffle surplus but cannot remove guarantees.
\end{qa}

\begin{qa}{549}{What is the environment the agent interacts with?}
\code{scheduler/env_seq.py}, a Gymnasium environment:
\begin{itemize}
\item a 5\,s step and 120-step episodes;
\item Markov context switches ($p_{stay}=0.92$);
\item AR(1) demand noise ($\phi=0.7$, $\sigma=0.15$);
\item random backhaul outages ($p=0.02$, 10 steps, core capacity halved);
\item work-conserving service with priority borrowing, backlog carry-over and fq\_codel-like 0.1\,s buffering;
\item the live allocator and constraint logic.
\end{itemize}
It is calibrated from 19 emulated runs (propagation floors 8/28\,ms, efficiency 1.0).
\end{qa}

\begin{qa}{550}{How is a state transition generated?}
\begin{enumerate}
\item Allocate (policy $\rightarrow$ priorities $\rightarrow$ allocator $\rightarrow$ constraints).
\item Serve: arrivals are demand times $\Delta t$ plus backlog; serve up to the allocation, lend the spare by priority.
\item Update backlog and overflow loss.
\item Advance context, demand noise and outage state.
\end{enumerate}
\end{qa}

\begin{qa}{551}{How are delayed rewards handled?}
Backlog carries consequences across steps; PPO uses $\gamma=0.95$ (about 20 steps, roughly 100\,s, of horizon) and
GAE advantages.
\end{qa}

\begin{qa}{552}{How do you handle a non-stationary environment?}
Non-stationarity (context switches, outages) is \emph{inside} the training distribution, and the observation exposes
context and backup state, so the policy can condition on it.
\honest{There is no online adaptation. A distribution shift beyond training (new contexts) is unhandled.}
\end{qa}

\begin{qa}{553}{Which baselines are needed to show that RL adds value?}
The rule-based demand-aware policy (identical pipeline, multipliers of 1) is the decisive one. Static, proportional and
3GPP baselines are also in the main study.
\end{qa}

\begin{qa}{554}{How do you evaluate an RL policy on scenarios it did not see in training?}
Simulator evaluation uses held-out seeds (\code{evaluate_policy}, seed 1000+). The live evaluation is in emulation,
which the simulator only approximates: a sim-to-emulation transfer test.
\honest{All contexts were seen in training. No unseen-context test.}
\end{qa}

\begin{qa}{555}{How do you prevent the agent exploiting simulator artefacts?}
Base-weight reward, the safety layer, calibration from emulation, and, decisively, evaluation in emulation (Mininet),
not just the simulator. The emulation shows no gain, consistent with no exploited artefact carrying over.
\end{qa}

\section{Training, safety and integration}

\begin{qa}{556}{Which algorithm: Q-learning, PPO, DQN, contextual bandits or something else?}
Two:
\begin{itemize}
\item the Cross-Entropy Method as a \emph{contextual bandit} (the original environment's actions do not affect
transitions, so this framing is exact);
\item PPO (stable-baselines3) on the sequential environment.
\end{itemize}
\end{qa}

\begin{qa}{557}{Why are those algorithms appropriate for the state and action spaces?}
Continuous low-dimensional actions with stochastic dynamics. PPO is a robust default for continuous control; CEM is a
strong derivative-free optimiser when the problem is a bandit. DQN would need a discretised action space.
\end{qa}

\begin{qa}{558}{Is the action space discrete or continuous?}
Continuous, 4-dimensional and bounded.
\end{qa}

\begin{qa}{559}{How is exploration controlled?}
PPO: a Gaussian policy with entropy bonus 0.005. CEM: the sampling distribution's standard deviation, refitted to the
elite 20\% each iteration (with a floor of 0.05).
\end{qa}

\begin{qa}{560}{What is the training sample complexity?}
PPO: 500k environment steps (8 parallel environments), minutes on CPU. CEM: 50 iterations $\times$ 32 samples per
context.
\end{qa}

\begin{qa}{561}{How do you know training has converged?}
\honest{Only through final-policy evaluation. Learning curves were not logged (SB3 verbose off). Add
TensorBoard or CSV logging and show the reward plateau over timesteps.}
\end{qa}

\begin{qa}{562}{How is the reward computed and logged?}
Environment: per step, reported per episode by \code{evaluate_policy} (mean, std, violations, satisfaction) in
\code{results/rl_training.json}. Live runs: per tick in \code{feedback.csv} (same reward definition, measured inputs).
\end{qa}

\begin{qa}{563}{Is the agent actually learning from reward, or fitting outcomes of random actions?}
\textbf{Today it learns from reward}: CEM keeps the elite by reward; PPO maximises advantage-weighted likelihood.
\textbf{The original implementation did not.} It fitted a regressor to \emph{uniformly random} actions and ignored
reward entirely, which is behaviour cloning of noise. This was found in the verification audit and replaced. Say this
openly: it shows the audit worked.
\end{qa}

\begin{qa}{564}{How do you distinguish a learned policy from a predictive model or regression heuristic?}
A learned policy improves its \emph{return} through interaction (CEM's elite reward rises over iterations; PPO's
episode reward against random and against rule). A regression heuristic predicts something without optimising
return. The old code was the second kind.
\end{qa}

\begin{qa}{565}{Is the trained agent wired into the live allocation or recovery loop?}
Allocation: yes. \code{rl_cem} and \code{rl_ppo} schedulers plug policies into the Priority Engine; they ran 25 main-study
runs each in emulation (L3). Recovery: no.
\end{qa}

\begin{qa}{566}{If it were not wired in, what would the RL implementation contribute?}
It is wired in. Its contribution is a negative empirical result plus a safe integration pattern (policy $\rightarrow$
priority multipliers $\rightarrow$ allocator with floors and constraints), and a calibrated environment for future
policies.
\end{qa}

\begin{qa}{567}{How would you integrate it without letting unverified actions modify the network?}
That is the current pattern: the policy can only rescale weights; floors and constraints still apply; outputs are
clipped; a missing model falls back to the rule (recorded). For more autonomy: shadow mode first, then an action
budget.
\end{qa}

\begin{qa}{568}{Offline training, online learning, or both?}
Offline in the calibrated simulator, with calibration from emulation feedback. Online learning in a critical network
would need safe exploration and is not done.
\end{qa}

\begin{qa}{569}{How do you constrain the action space to preserve minimum QoS guarantees?}
Bounded multipliers; the floors and constraints sit \emph{outside} the policy, so no action can violate them.
\end{qa}

\begin{qa}{570}{How would you compare RL with static priority, greedy allocation and the recovery planner?}
Done for allocation in the main study (7 schedulers, 25 paired runs each, Wilcoxon). Not done for recovery (RL does
not act on recovery).
\end{qa}

\begin{qa}{571}{How do you handle the absence of GPU infrastructure without unsupported claims about training requirements?}
Everything was trained on CPU: PPO (500k steps) in minutes, SliceNet-CNN (\NCnnParams{} parameters) in under a minute.
We claim nothing that needs a GPU.
\end{qa}

\begin{qa}{572}{Which experiments can run on CPU, and which genuinely need acceleration?}
All of ours run on CPU. Accelerators would matter for much larger models (Transformers over raw packet sequences),
large-scale hyper-parameter search, or simulators with thousands of parallel environments.
\end{qa}

\begin{qa}{573}{What is the safest way to introduce a learned policy into a live network?}
Shadow mode (log the would-be actions); a constrained action space and safety layer (already present); a canary site;
automatic rollback on SLA regression; a human-in-the-loop approval for policy updates.
\end{qa}

\begin{qa}{574}{What evidence would justify keeping RL in the final architecture rather than as future work?}
A significant, replicated improvement in emulation (paired across seeds and contexts, CI excluding zero) on the
critical-service metrics, without a fairness regression. We do not have it, so RL is future work, and the
environment plus integration pattern are the reusable artefacts.
\end{qa}
